\section{From points on determinant rings to the two prime conditions}

Let $T$ be an original setup, $n=2$ or $3$, $k=O/\mathfrak m_O$,
$B_d$ its universal determinant ring, $B\subset A$ the coefficient
and entry algebras of its matrix space. For any comparison map
$\Phi:B_d\to B$ satisfying $\mathrm{IsPhi}$ and original chart
$L$ write $R=L.R$, $b:B_d\to B\to A\to R$, and
$q(P)=b^{-1}P$. Write $G$ for the arithmetic group $T.G.G_0$,
$G'=\ker\varepsilon$ within it, $\pi:G\twoheadrightarrow F.G_0$
for the link projection and $\kappa$ for its scalar character.
On a specialized matrix representation $r$, the arithmetic
representation is $g\mapsto\kappa(g)r(\pi g)$, with local
characters twisted by the same scalar at the linked local element.
In this chapter $J_d$ is an ideal of $B_d$.

% BEGIN P68K1
\paragraph{P68K1. The power-series point with prescribed kernel.}
Suppose $p\subset R$ is prime, $\dim R/p=1$, and $\ell\in p$.
The contraction of $p$ to the coefficient DVR contains the nonzero
element $\ell$ by $T.\mathrm{lne}$, hence is $\mathfrak m_O$.
The complete Noetherian local domain $R/p$ has characteristic
$\ell$ and finite residue field by the local chart axioms. Its
normalization in its fraction field is finite (a complete
Noetherian local ring is excellent), and is local: it is a
domain finite over a Henselian local ring, hence the product
decomposition into local factors has just one factor. The
normalization is therefore a complete one-dimensional normal
Noetherian local domain, a DVR. Its residue field $K$ is finite.
In equal characteristic, the simple roots of $X^{|K|}-X$ lift
uniquely by Hensel. These lifts form a field by the Frobenius
identity. Expanding along a uniformizer identifies the DVR with
$K[[t]]$. This gives a local map $j:R\to K[[t]]$ of kernel $p$.
The image of $k$ is constant: its elements satisfy
$X^{|k|}=X$, whose roots in the power-series ring are the
unique lifts of their constant terms. Give $K$ this induced
$k$-algebra structure. Then $j$ is an $O$-algebra map.

Use the linked twist of $j_*L.\rho$ for the arithmetic matrices and
characters, with symmetric antidiagonal form $J_n$. The star
identity for $L.\rho$ gives the form equation because
$\kappa(g)\kappa(g^c)=\varepsilon(g)^{1-n}$.
Continuity follows from ALC17, the continuity of $\pi,\kappa$
and locality of $j$. The comparison identity extends along $jb$:
by injectivity of $A$ in the complete matrix ring it already holds
there before applying $j$, as for the polynomial laws in W67F3.
The map $B_d\to K$ given by the constant term of $jb$ has finite
domain image, hence maximal kernel. Since $B_d$ is local with
residue $k$, this map is its specified residue map followed by
$k\to K$. It gives the required residual determinant and each
ordered residual character. At a local separator, the residual ordinary characters are the distinct
powers of $\Psi$, and ACP4 gives their product characteristic
polynomial. The normalized ordered relation and ACP3 therefore give
the ordered local flag; it extends to the fraction field and persists
under a scalar twist, with all its diagonal characters twisted.

For $g\in I_p(u)$ every residual character equals one. The linked
tameness data give $m$ prime to $\ell$ with $\pi(g)^m=\kappa(g)^m=1$.
Thus the twisted matrix $Z$ satisfies $Z^m=1$ and its residual
characteristic polynomial is $(X-1)^n$. The sum
$1+Z+\cdots+Z^{m-1}$ is a unit matrix: modulo the maximal ideal
its eigenvalues are all $m\ne0$. Multiplying it by $Z-1$ gives
zero, so $Z=1$. These checks produce a point $y$ in the sense
of the original determinant problem. It has $y.A=K[[t]]$,
the exact twisted matrices and characters, and $\mathrm{IsPt}(y,jb)$.
This proves $T.\mathrm{PointEx}$ without requiring $\Phi$ surjective.
% END P68K1

% BEGIN P68K2
\paragraph{P68K2. A local element separating all positive powers.}
For every point and local $j$ with the matching twisted matrices
and characters, suppose $\ker j$ is nice on the chart. At a fixed
place $v$ let $N$ be the inverse image of $F.N_{e(v)}$ under the
local group isomorphism. Its image in $G$ belongs to $G'$.
The ordinary ratios $\chi_{ij}=y.\psi_i/y.\psi_j$ on $N$
reduce to one, since the residual $\Psi$ is one there and the
common twist cancels. Each is nontrivial by the noncoincidence
clause of niceness on this exact subgroup. The group
$1+tK[[t]]$ is torsion-free: an $\ell$-power acts by Frobenius,
and a power prime to $\ell$ preserves the order of the first
nonzero coefficient of $u-1$. In particular each ratio has
infinite image. For any open subgroup $U\subset G$, its
local inverse image has finite index, so all ratios are still
nontrivial on its intersection with $N$.

One element works simultaneously. In the joint image under the
ratios the group law is abelian. For finitely many nontrivial
characters into a torsion-free group, start with an element on
which the already treated characters are nontrivial. If the next
character is trivial on it, multiply it by $h^a$ for an $h$
on which the new character is nontrivial. For each previously
treated character at most one exponent is prohibited if its
value on $h$ is nontrivial, and none otherwise. A positive
exponent outside this finite list gives the induction. We obtain
a local element with image in $U\cap G'$ whose ordered character
values have non-torsion pairwise ratios. Its characteristic
polynomial splits with those eigenvalues by the flag. With
$U=G$ this gives every distinctness and cyclotomic-twisted
distinctness in $y.N1$, since $\varepsilon=1$ on $G'$.
% END P68K2

% BEGIN P68K3
\paragraph{P68K3. The irreducible constituents on the cyclotomic kernel.}
Put $K_*=\overline{\operatorname{Frac}(y.A)}$. In type I the
full representation over $K_*$ is irreducible by ALC15 and
ALC17: the kernel on the entry ring is the contraction of
$\ker j$, and invariant subspaces are unaffected by the scalar
twist and the surjection $\pi$. In type E the odd entries vanish.
The two outer entries $02,20$ each survive somewhere because
the point is off triangular. Projectors onto the normalized
coordinates are linear combinations of universal group matrices
over the coefficient ring by ALC13. On the outer block,
every invariant subspace is thus coordinate-spanned and the
surviving entries in both directions make that block irreducible.

For either of these irreducible representations $V$, its dimension
$r$ is prime. We show irreducibility on $H=U^{\rm nor}\cap G'$
for every open $U\subset G$, where $U^{\rm nor}$ is its open
normal core. A local place exists since $F.\mathrm{Big}$ gives
$d>0$. On its local group $D$ the characteristic polynomial on
$V$ is the product of the indicated local characters. For the
outer block, the middle line equals one fixed local character:
in the middle entry of the ordinary ordered product, one may
choose the group argument independently in each factor. If no
one character agreed identically, one could make every factor
nonzero. Cancel this middle factor from the full characteristic
polynomial to obtain the other two local characters. P68K2
separates all these characters on $D\cap H$ (inverse image in $D$).

There are integral characteristic polynomials on $V$ with residual
roots given by the global residual characters: this is the
residual determinant in the full case. For the outer block, ALC17
gives residual diagonal values. The map from $B$ to the finite
residue field of the chart has finite field image, so its kernel
is the maximal ideal. ALC12 therefore makes the residual matrix
triangular. Thus its residual roots are the two outer
characters. In both cases some element of $D$ separates two
global residual characters, since the frame supplies
$\Psi(s)^2\ne1$ and a common scalar twist preserves inequality.

The sum of the $G$-translates of any simple $H$-submodule is $V$
by normality and irreducibility. Hence the restriction is
semisimple, and $G$ transitively permutes its isotypic blocks.
Their number, common multiplicity and common simple dimension
multiply to the prime $r$. Besides irreducibility the two
possibilities are a scalar action by $H$, or a transitive system
of $r$ distinct lines. A local element of $H$ with simple
spectrum excludes the scalar case. In the line case let $N_0$
be the kernel of the resulting permutation action of $G$.
On $N_0$ the representation is the sum of line characters
$c_i$, and $H\subseteq N_0$. On $D\cap H$ the multiset of these
characters is the multiset of the local characters, by equality
of the split characteristic polynomials. Here and below this
multiset assertion works in positive characteristic as well:
an integral relation between character multisets which vanishes
pointwise as a formal multiset gives a linear relation between
the distinct multiplicative characters. Their linear independence
forces every multiplicity difference divisible by the characteristic;
dividing the original integral multiset identity and repeating
forces the differences to vanish. In characteristic zero a single
application of independence suffices.
The local characters are invariant under conjugation by $D$
and distinct on $D\cap H$. Consequently $D$ fixes each line
and lies in $N_0$.

Extend the residue of $y.A$ to the integral closure in $K_*$
by choosing a prime above its maximal ideal, and embed the
resulting residue field in an algebraic closure of the residual
field. Both $c_i(g)$ and their inverses, for $g\in N_0$, are
integral roots, so yield reduced characters. The same multiset
argument on $N_0$ identifies them up to permutation with the
global residual characters restricted there. Conjugation permutes
the $c_i$ transitively and each global residual character is
conjugation-invariant. Thus all reduced characters agree on
$N_0$, in particular on $D$, contrary to the residual separator.
This proves irreducibility on $H$, hence on $U\cap G'$.
In particular type I gives $y.\mathrm{TypeI}$. In type E, reorder
the coordinates $(0,2;1)$ to obtain the original $2+1$
decomposition over $\operatorname{Frac}(y.A)$, with forms the
antidiagonal $J_2$ and $1$, both with the original weight
$1-3$. Both blocks are absolutely irreducible on $G'$ and
any equivariant map from the plane to the line vanishes, since
a nonzero such map would have an invariant line as kernel.
These are all clauses of $y.\mathrm{TypeE}$.
% END P68K3

% BEGIN P68K4
\paragraph{P68K4. Regular elements in every required open set.}
Take an open set meeting $G'$, choose $s$ in the intersection
and an open subgroup $V\subset G$ whose translate $sV$ lies
in the set. Over $K_*$ let $H$ be the Zariski closure of the
image of $V\cap G'$. Its identity component $H^0$ acts
irreducibly on the full space in type I and on the outer plane
in type E. Indeed the inverse image of $H^0$ is open in
$V\cap G'$: each of the finitely many components is cut out
there by polynomial equations, whose zero sets are closed for
the coefficient valuation. The original matrices are continuous
for that valuation, and the finitely many coefficients of each
polynomial belong to a finite extension of the Laurent-series
field. A sufficiently small open subgroup $V'$ of $G$ therefore
has $V'\cap G'$ in this inverse image, and P68K3 applies.

By P68K2 the group contains an element $u$ with pairwise
non-torsion eigenvalue ratios. In an eigenbasis the identity
component $S$ of the Zariski closure of its powers is a split
torus in $H^0$. The diagonal characters of $S$ are distinct:
otherwise the quotient of two eigenvalues would have finite
order on the component group. Choose an integral cocharacter
of $S$ outside the finitely many equality hyperplanes. It has
pairwise distinct integral weights $a_1<\cdots<a_n$, in reordered
eigenbasis. In type E use an eigenbasis adapted to the plane
and line, which is possible by the simple spectrum.

There is $h\in H^0$ for which the leading principal minors
of $Y=y.\rho(s)h$ in this ordering are all nonzero. For each
size $i$, this minor is the corresponding matrix coefficient on
$\bigwedge^i$. In the irreducible case $n\le3$, the nontrivial
exterior powers are the standard representation and its dual
times the determinant (or just the standard one). They are
irreducible on $H^0$. The orbit of the coordinate wedge thus
spans, and the nonzero row functional after $y.\rho(s)$ is nonzero
on some orbit element. In the $2+1$ case the same coordinate
wedge lies in a pure line, the outer plane, or that plane tensored
with the middle line. Each such space is irreducible and is
preserved by $y.\rho(s)$, so the same matrix coefficient is
again nonzero somewhere. Each minor supplies a nonempty Zariski
open of the irreducible group $H^0$; their intersection is nonempty.

Multiply $Y$ on the right by
$\operatorname{diag}(z^{a_1},\ldots,z^{a_n})$. Up to sign the
coefficient of $X^{n-i}$ in its characteristic polynomial is
the sum over size-$i$ principal minors with monomials
$z^{\sum_{j\in I}a_j}$. The unique least exponent is
$a_1+\cdots+a_i$ and its coefficient is nonzero by construction.
Thus the Newton polygon in the $z$-adic valuation has $n$
distinct slopes, each of horizontal length one. The roots have
distinct valuations and the discriminant is nonzero. Some
$z\in K_*^\times$ keeps it nonzero. This proves that the regular
open in $y.\rho(s)H$ is nonempty. By density there exists an
actual $h'\in V\cap G'$ with separable characteristic polynomial
at $s h'$. Separability is unchanged by the field extension to
$K_*$. This proves $y.N3$ in both cases, including when the
original open set does not contain the identity. P68K2--P68K4
give the full native $T.\mathrm{NiceOf}$.
% END P68K4

% BEGIN P68K5
\paragraph{P68K5. A type-I prime gives a stable determinant prime.}
For any prime $P\subset R$ of type I let
$K'=\overline{\operatorname{Frac}(R/P)}$ and
$k_q=\overline{\operatorname{Frac}(B_d/q(P))}$. Embed $k_q$
in $K'$ over $B_d/q(P)$. By ALC15 and ALC17 the underlying
entry representation in $K'$ is irreducible, hence the same
holds after pullback by $\pi$ and twist by $\kappa$. Its
determinant law comes from $k_q$ by $\mathrm{IsPhi}$.
We verify descent of irreducibility to the field in the definition
of stability. By Burnside, choose $n^2$ group matrices $M_j$
forming a basis of the full matrix algebra in $K'$. All traces
of words lie in $k_q$ because the determinant law does. The
trace pairing on the full matrix algebra is nondegenerate in
every characteristic, by the matrix units. Thus the Gram matrix
$(\operatorname{tr}(M_iM_j))$ is invertible over $k_q$. Solving
with this matrix puts the structure constants for $M_iM_j$,
the identity, and the coordinates of every group matrix in
$k_q$. Their $k_q$-span is a central simple algebra: extension
to $K'$ gives a matrix algebra, so proper two-sided ideals and
extra center are excluded by faithful scalar extension. It
splits over the algebraically closed field $k_q$. The resulting
$n$-dimensional representation is irreducible because the group
matrices span its matrix algebra. After extension it is conjugate
to the original one by the innerness of automorphisms of matrix
algebras. Its determinant law therefore agrees over $k_q$ by
injectivity of scalar extension on polynomial coefficients. This
is precisely stability of $q(P)$, proving $T.\mathrm{StableOf}$.
% END P68K5

% BEGIN P68K6
\paragraph{P68K6. A thick prime gives a nonsemisimple determinant algebra.}
In rank three suppose $P\subset R$ is thick. Put $Q=R/P$,
$f:B_d\to Q$, and choose its native thick representation $s_0$
over $Q$. Its even entries equal those of the block-diagonal
ordinary matrices, its middle row vanishes off the diagonal,
and it satisfies both ordered products, including the reversed
star product. Over every embedding of $Q$ in a field the outer
plane has no equivariant section of its one-dimensional quotient.
Use $s(g)=\kappa(g)s_0(\pi g)$. The determinant law of $s$
is $f_*D$: for any scalar extension and any linear combination
of group matrices, $s$ is block triangular and has the same
two diagonal blocks as the split specialization. The two
determinants are equal. We show that its group-algebra action
annihilates the actual closed ideal defining the determinant
algebra.

The Cayley--Hamilton generators vanish by the matrix identity,
including the coefficients obtained after polynomial base change.
An ordered local generator maps to the product of twisting
scalars times the original zero product. On a group term the
star sends $a[g]$ to
$a\varepsilon(g)^{1-3}[(g^c)^{-1}]$. Its image under $s$ is
$f(a)\kappa(g)s_0((\pi g)^*)$ by the twist identity. Since star
reverses products, the other local generator vanishes by the
given reversed star equation for $s_0$. On $I_p(u)$ tameness
gives an invertible prime-to-$\ell$ exponent killing both the
linked element and its scalar. The residual characteristic
polynomial is $(X-1)^3$, using the residual determinant as in
P68K1. The coefficient map $B_d\to Q$ is local because its
reduction to the finite residue field has a finite field as image,
hence maximal kernel. The invertible geometric-sum argument of
P68K1 applies also over $Q$, so $s(g)=1$. It applies to the
star term as well: its twisted-star image has prime-to-$\ell$
power one by the same tame conditions on $\pi g$ and $\kappa(g)$.
Its characteristic polynomial equals that of the unstarred
element by the polarization clause of the universal determinant,
which equates every signed coefficient for these two group-algebra
elements. In particular its residual eigenvalues are all one. Thus both inertia generators vanish.

Modulo each $\mathfrak m_Q^N$ the entries of $s$ are locally
constant. Their representation has open normal kernel. Reducing
the group algebra by this kernel and $\mathfrak m_{B_d}^N$ gives
one of the quotients in the definition of the closed ideal.
Together with generator vanishing this proves its image lies in
every $\mathfrak m_Q^N$ entrywise, hence is zero by Krull
intersection. Consequently, with $k_q\subset K'$ as in P68K5,
the $K'$-space carrying $s$ is a module over the source's
$E^0_{k_q}$, viewed as a ring over the smaller field. Denote
its middle quotient character by $m$. Each $m(g)$ belongs to
$k_q$: it is a root of the characteristic polynomial over the
algebraically closed $k_q$ supplied by the determinant law.

The prime $q(P)$ cannot be stable. If an irreducible representation
over $k_q$ gave the same determinant, its group algebra would
surject onto $M_3(k_q)$ by Burnside. On an element in the kernel
of this matrix action the common characteristic polynomial is
$X^3$, so its action on the middle quotient has eigenvalue zero.
The middle character on the group algebra would thus factor as
a unital map $M_3(k_q)\to K'$, impossible by the matrix-unit
relations. Nor can $E^0_{k_q}$ be semisimple: its module $K'^3$
would then have a module complement to the submodule of vectors
with middle coordinate zero. Choose a vector $v$ in that
complement with middle coordinate one by decomposing the middle
basis vector. For every $g$, both $s(g)v$ and $m(g)v$ belong
to the complement; the latter assertion uses $m(g)\in k_q$.
Their difference has middle coordinate zero and must vanish.
Cancel $\kappa(g)$ and use surjectivity of $\pi$ to obtain
the forbidden section for $s_0$. We have proved both negations
defining $T.B.\mathrm{IsThick}(q(P))$, hence $T.\mathrm{ThickOf}$.
% END P68K6

% BEGIN P68K7
\paragraph{P68K7. Transport of the two point conditions, including trivial chains.}
Write $M=J_dB$. For any chart ideal $H$,
$M\subseteq\pi_L(H)$ if and only if $J_d\subseteq b^{-1}(H)$.
Given a good nice type-I prime $p$, P68K1 realizes it as
$\ker j$ and supplies $\mathrm{IsPt}(y,jb)$. P68K4 and the
given $N4$ give $y.\mathrm{NiceI}$. Goodness means that $jb$
kills $J_d$. For each prime $P\subseteq p$, $q(P)$ lies under
$\ker(jb)$. The given $P1_{\rm core}$ forces $J_d\subseteq q(P)$,
which is precisely goodness of $P$. This proves
$T.\mathrm{FirstPrimeConditionTransfer}$ in both ranks.

For $n=3$ enumerate a source thick-link chain as
$p_0,\ldots,p_h$, with thick prime $E_i$ contained in
$p_i,p_{i+1}$ for $0\le i<h$. Each $p_i$ is nice type E,
including $p_0$ when $h=0$. Choose P68K1 points
$(y_i,j_i)$ individually on the given chart; each has
$\mathrm{IsPt}(y_i,j_ib)$ and is $\mathrm{NiceE}$ by P68K4
and the supplied $N4$. Every $q(E_i)$ is thick by P68K6 and
belongs below both consecutive kernels. The supplied good
type-I prime $P\subseteq p_0$ yields a stable $q(P)$ by P68K5,
with $J_d\subseteq q(P)\subseteq\ker(j_0b)$. For any
type-I $P'\subseteq p_h$, $q(P')$ also is stable and lies
below $\ker(j_hb)$. Apply the original $P3_{\rm core}$ to
these $h+1$ points, $h$ thick primes, and the indicated initial
and final stable primes. It gives $J_d\subseteq q(P')$, exactly
the required goodness. At $h=0$ the middle list is empty and
the same argument applies. This proves
$T.\mathrm{SecondThirdPrimeConditionTransfer}$ and hence supplies
$P1'$ and $P23'$ for the original image ideals from the appropriate
$\mathrm{ModIn}$ fields.
% END P68K7
